% Independent linear reconstruction audit. Input fragment, not a standalone file.
\section{A regular particular metric from the sourced conformal equations}
\label{sec:sourced-cefe-reconstruction}

This section concerns the \emph{linearized} problem only.  It supplies a
particular physical metric, before imposing the final elliptic shift and
lapse gauges.  It assumes the strong subsidiary closure and the compatible
regular rescaled Weyl evolution established in the preceding sections.
It does not infer regularity of the rescaled Weyl tensor from the subsidiary
weights alone.  Initial Weyl peeling and curvature compatibility are explicit
parts of the initial-data class.

Write $g$ for the conformal metric in this section, $\widetilde g=\Omega^{-2}g$
for the physical metric, $\Sigma_a=\nabla_a\Omega$, and $L_{ab}$ for the
\emph{unphysical} Schouten tensor.  The prescribed first-order physical Ricci
source is
\[
 \delta\widetilde R_{ab}=\Omega^2r_{ab},\qquad
 \ell_{ab}=\tfrac12(r_{ab}-\tfrac16g_{ab}\operatorname{tr}_g r),
 \qquad \delta\widetilde L_{ab}=\Omega^2\ell_{ab}.
\]
Every coefficient multiplying a first-order source in this section is evaluated
on the fixed conformal Minkowski background.  Thus changes of coordinates or
conformal gauge do not change $r,\ell,J,d$ at first order: the physical Ricci
and Weyl background fields vanish.

\subsection{The source conservation identity, including its trace}

Put $S_{ab}=r_{ab}-\tfrac12g_{ab}\operatorname{tr}r$.  The connection-change
formula gives, without using a field equation,
\[
 \widetilde\nabla^a(\Omega^2S_{ab})
 =\Omega^4\nabla^aS_{ab}+\Omega^3\Sigma_b\operatorname{tr}S.
\]
Consequently the necessary and sufficient linear physical conservation law is
\begin{equation}
 \Omega\bigl(\nabla^a\ell_{ab}-\nabla_b\operatorname{tr}\ell\bigr)
 =3\Sigma_b\operatorname{tr}\ell.
 \label{eq:sourced-cefe-conservation}
\end{equation}
It is not legitimate to prescribe arbitrary smooth $r$ independently of the
subsidiary system.  Here the $H,M$ evolution equations were obtained from the
physical ADM contracted Bianchi identity with the stated $Q^{\rm phys}_{ij}$.
Substituting their Ricci projections into that identity is precisely
\eqref{eq:sourced-cefe-conservation}; it introduces no new condition on their
solutions.  This explains why the source equations must be kept, rather than
regarding $H,M,T,z$ as unrelated external functions.
For a direct normalization check set $N=A/\Omega$ and
$F_{ij}=Q^{\rm phys}_{ij}/N$.  The prescribed Einstein source has projections
\[
 \delta\widetilde G_{nn}=H/2,\qquad
 \delta\widetilde G_{ni}=-M_i,\qquad
 \delta\widetilde G_{ij}
  =\gamma^{\rm phys}_{ij}(H/2+\operatorname{tr}_{\rm phys}F)-F_{ij}.
\]
Indeed $\widetilde R_{ij}=-F_{ij}$ and
$\widetilde R_{nn}=H+\operatorname{tr}_{\rm phys}F$.
Its normal and tangential divergence projections are exactly the displayed
closed Hamiltonian and momentum subsidiary equations, including the
covector Lie derivative on the latter.  This provides an explicit algebraic
way to verify conservation rather than assuming a physical metric has
already been reconstructed.

The identity also has a nontrivial boundary consequence.  Since $d\Omega\ne0$
at scri and the source is regular, $\operatorname{tr}\ell=O(\Omega)$.  For the
explicit source in the Weyl section,
\[
 \operatorname{tr}r=-h_4-12G\cdot\zeta
       -4\Omega\partial_i\zeta^i+6\kappa\tau_4.
\]
The root-two subsidiary coefficient relation makes this trace vanish at scri.
This is consistent with the full identity, but checking this one leading
coefficient alone would not prove conservation.

The conformal Cotton tensor is the already established regular source
\[
 J_{bcd}=\Omega(\nabla_c\ell_{db}-\nabla_d\ell_{cb})
 +3(\Sigma_c\ell_{db}-\Sigma_d\ell_{cb})
 -g_{cb}\Sigma^e\ell_{de}+g_{db}\Sigma^e\ell_{ce}.
\]
For example its contracted trace vanishes exactly by
\eqref{eq:sourced-cefe-conservation}.  Its remaining integrability identities
follow by commuting the physical derivatives on the flat physical background.

\subsection{Exact regular equations for a traceful source}

The usual trace-free-matter conformal theorem cannot simply be cited for this
source.  Introduce instead
\[
 s=\frac{\Sigma_a\Sigma^a}{2\Omega}
   =s_F-\frac14\Omega^3\operatorname{tr}\ell,
 \qquad s_F=\tfrac14\Box_g\Omega+\tfrac1{24}R[g]\Omega.
\]
The quotient is an identity on $\Omega>0$; in the reduced system $s$ is an
independent regular field.  The exact sourced conformal equations are
\begin{align}
 H_{ab}:=\nabla_a\nabla_b\Omega
  &=-\Omega L_{ab}+sg_{ab}+\Omega^3\ell_{ab},
       \label{eq:sourced-cefe-hessian}\\
 \nabla_as&=-L_{ab}\Sigma^b+\Omega^2\ell_{ab}\Sigma^b,
       \label{eq:sourced-cefe-s}\\
 \nabla_cL_{ab}-\nabla_aL_{cb}
  &=\Sigma_ed^e{}_{bca}+\Omega J_{bca}=:F_{cab},
       \label{eq:sourced-cefe-curl}\\
 \nabla_ad^a{}_{bcd}&=J_{bcd},
       \label{eq:sourced-cefe-bianchi}\\
 R_{abcd}[g]&=\Omega d_{abcd}
    +2\bigl(g_{a[c}L_{d]b}-g_{b[c}L_{d]a}\bigr),
       \label{eq:sourced-cefe-curvature}\\
 2\Omega s-\Sigma^2&=0.
       \label{eq:sourced-cefe-scalar}
\end{align}
Antisymmetrization in the curvature formula has weight $1/2$.
The Schouten conformal transformation proves
\eqref{eq:sourced-cefe-hessian}; differentiation of
\eqref{eq:sourced-cefe-scalar} then proves
\eqref{eq:sourced-cefe-s}.  The Weyl conformal transformation and the two
contracted Bianchi identities prove the curl and divergence equations.
In particular all new terms are $\Omega^3\ell$, $\Omega^2\ell\Sigma$,
and $\Omega J$, with no inverse conformal factor.

\subsection{An explicit regular wave reduction}

Prescribe the smooth unphysical scalar curvature $\mathcal R(x)$ and smooth
generalized harmonic coordinate sources.  Impose $R[g]=\mathcal R$ and use
the usual reduced Ricci operator, so its principal part is
$-\tfrac12g^{cd}\partial_c\partial_dg_{ab}$.
The following equations define a regular reduced system.  In their right-hand
sides replace every curvature by the algebraic expression
\eqref{eq:sourced-cefe-curvature} and every Hessian of $\Omega$ by
\eqref{eq:sourced-cefe-hessian}:
\begin{align}
 R^{\rm red}_{ab}[g]&=2L_{ab}+\tfrac16\mathcal R g_{ab},
        \label{eq:sourced-reduced-metric}\\
 \Box_g L_{ab}&=\tfrac16\nabla_a\nabla_b\mathcal R
   +\nabla^cF_{c(ab)}
   +[\nabla^c,\nabla_{(a}]L_{|c|b)},
        \label{eq:sourced-reduced-L}\\
 \Box_g\Omega&=4s-\tfrac16\mathcal R\Omega
                    +\Omega^3\operatorname{tr}\ell,
        \label{eq:sourced-reduced-Omega}\\
 \Box_gs&=-\tfrac16\Sigma^a\nabla_a\mathcal R-L^{ab}H_{ab}
              +\nabla^a(\Omega^2\ell_{ab}\Sigma^b).
        \label{eq:sourced-reduced-s}
\end{align}
The symmetric $L$ equation follows by taking the divergence of its curl
equation and using $\nabla^aL_{ab}=\nabla_b\mathcal R/6$.
More explicitly,
\[
 \nabla^cF_{cab}
 =H^c{}_e d^e{}_{bca}+\Sigma_e\nabla^cd^e{}_{bca}
    +\Sigma^cJ_{bca}+\Omega\nabla^cJ_{bca}.
\]
Thus the largest source derivatives are first derivatives of $d,J,\ell$.
This formula also specifies the reduction without concealing a second
derivative of $\Omega$ in $\nabla F$.  All these statements concern the
linearization of the displayed system about the fixed background; no
quasilinear conclusion is needed here.

In coordinates the scalar wave operators must be reduced using the same
contracted-Christoffel gauge.  This matters for the mixed regularity below:
$\delta\Box\Omega_*$ then contains only zeroth-order metric perturbations
and prescribed gauge terms.  Leaving an uncontrolled contracted
$\delta\Gamma^a\partial_a\Omega_*$ on the right would falsely remove the
one extra derivative needed for the final conformal-factor gauge change.

\subsection{Mixed estimate and the fixed-$\Omega$ particular metric}

For integer $m\geq3$ use the phase norm
\begin{align*}
 \mathcal E_m={}&\|\delta g\|_{H^{m+1}}+
 \|\partial_t\delta g\|_{H^m}
 +\|\delta L\|_{H^m}+\|\partial_t\delta L\|_{H^{m-1}}\\
 &+\|\omega\|_{H^{m+2}}+\|\partial_t\omega\|_{H^{m+1}}
 +\|\delta s\|_{H^{m+1}}+\|\partial_t\delta s\|_{H^m},
 \qquad\omega=\delta\Omega.
\end{align*}
One sufficient source norm is
\begin{align*}
 \mathcal S_m={}&\|d\|_{H^m}+\|\partial_td\|_{H^{m-1}}
 +\|J\|_{H^m}+\|\partial_tJ\|_{H^{m-1}}\\
 &+\|\ell\|_{H^{m+1}}+\|\partial_t\ell\|_{H^m}.
\end{align*}
The independent Bianchi evolution bounds the displayed time derivative of
$d$ by its spatial derivatives and $J$.
The ordinary energy estimate for the linear reduced waves gives
\begin{equation}
 \mathcal E_m(t)\leq C_{m,T}\left(
 \mathcal E_m(0)+\int_0^t\mathcal S_m(t')\,dt'\right),\qquad 0\leq t\leq T.
 \label{eq:sourced-cefe-mixed-energy}
\end{equation}
To verify the derivative count, the $L$ forcing belongs to $H^{m-1}$;
even the tensor wave operator's second metric derivatives have this order
when $g\in H^{m+1}$.  The metric forcing belongs to $H^m$, the $s$ forcing
to $H^m$, and the $\Omega$ forcing to $H^{m+1}$.  The last assertion uses
the contracted-Christoffel gauge just specified.  The background light cone
is the common wave principal symbol.  At scri $B^n=-A$, so all reduced wave
fluxes are nonpositive.  Standard extension from the ball, or the equivalent
outflow energy construction, gives existence and uniqueness without an
independent outer boundary value.  Commuting finitely many time derivatives
and rotations gives the same estimate with the corresponding joint source
norms.  No norm count is asserted for an infinite commutator hierarchy.

Finite conormal commutations require a separate argument.  In a collar
write each regular first-order wave reduction as
$\mathcal A=V^r\partial_r+V^A\partial_A+R$ and set
$\mathcal D=b(r)\partial_r$, $b=\Omega/\Omega_r$, extended by a cutoff
away from the centre.  All matrices here are smooth and bounded.  The exact
identity is
\[
 [\mathcal A,\mathcal D]=b'\mathcal A
 -(b'V^A+b\partial_rV^A)\partial_A
 -(\partial_rV^r)\mathcal D-b'R-b\partial_rR.
\]
On a solution replace $\mathcal AU$ by $U_t$ minus its known source.
In the equation for $\mathcal D^jU$, the time and angular derivatives then
have strictly smaller conormal order and the same total joint order; the
term involving $\mathcal D(\mathcal D^{j-1}U)$ is bounded multiplication at
the current level.  Rotations span the angular derivatives on the collar.
Induction on conormal order at fixed total order and the same energy estimate
therefore close the joint time/rotation/conormal hierarchy.  The first-order
gradient compatibilities recover the stated ordinary Sobolev orders.
The $\nabla d$ term in the $L$ equation remains an upstream source
($d\in H^m$ gives $\nabla d\in H^{m-1}$); the other mixed-phase couplings
have the bounded orders already checked above.  Cutoff commutators are
supported in the interior and use ordinary interior derivatives.
Include three cutoff Cartesian generators $\chi\partial_i$, equal to the
Cartesian derivatives near the centre and supported away from scri, in the
joint family $\mathcal Z$ along with $\mathcal D$ and rotations.  The cutoffs
overlap so the family spans ordinary derivatives uniformly in the interior.
Their commutators are supported where $\Omega$ is bounded below and are the
ordinary bounded-coefficient differentiated wave system.  Thus the explicit
joint phase is
\[
 \mathcal E_{4,32}
   =\sum_{p+|I|\le32}\mathcal E_4
           (\mathcal Z^I\partial_t^p(\delta g,\delta L,\omega,\delta s)).
\]
The compatible derivative fields and their velocities are interpreted as in
the mixed phase above.  This norm includes all corresponding ordinary
interior derivatives, in particular the metric/lapse interior orders needed
by the gauge-wave source.  The source phase uses the same generator words.
Consequently this finite joint hierarchy is controlled with its listed
source and initial norms.  This trades
conormal order for time/rotation order and does not assert extra ordinary
normal differentiability across a fractional branch.

To return to the original, fixed conformal factor, choose a smooth cutoff
$\chi$ supported in a collar, equal to one near scri, and a smooth spatial
vector $N$ on its support with $N(\Omega)=1$.  Set
\[
 \xi=-\chi\omega N,\qquad
 \vartheta=-\frac{(1-\chi)\omega}{\Omega}.
\]
The quotient is supported where $\Omega$ is bounded away from zero.  The
coordinate and conformal transformation
\[
 \omega'=\omega+\xi(\Omega)+\vartheta\Omega=0,\qquad
 \delta g'=\delta g+\mathcal L_\xi g_*+2\vartheta g_*
\]
preserves the physical metric modulo a regular coordinate transformation.
It is bounded from $\mathcal E_m$ to
$\delta g'\in H^{m+1}$, $\partial_t\delta g'\in H^m$.
The extra derivative of $\omega$ is precisely what controls the Lie
derivative.  At $m=4$ this yields a fixed-$\Omega$ particular conformal metric
$g_P\in H^5$, with its extrinsic-curvature perturbation in $H^4$.
The map to $h_P,c_P,u_P,v_P,a_P,k_P$ is algebraic in $g_P$ and its first
derivatives, hence has these same stated metric/curvature orders.  Finally
$q_P=k_P-2T$ and $\Lambda_P=V(h_P)+2z$ reconstruct the independently evolved
Z4c fields; those two fields must not be discarded during reconstruction.

This reconstruction recovers the full off-constraint geometry, rather than
projecting to vacuum constraints.  For any physical metric the contracted
Gauss and Codazzi identities give
\[
 H_{\rm geom}=\widetilde R_{nn}
       +\gamma_{\rm phys}^{ij}\widetilde R_{ij},\qquad
 M_{{\rm geom},i}=-\widetilde R_{ni}.
\]
The prescribed source therefore makes these exactly the independently
propagated $H,M$.  Since $\widetilde R_{ij}=-Q^{\rm phys}_{ij}/N$, its ADM
extrinsic-curvature evolution is precisely the tensor $C_{\rm Z4c}=0$
evolution with the stated sources.  The geometric metric evolution is the
ADM kinematic identity.  The independent $T$ equation is unchanged, and
the independent $z$ equation, with $\Lambda=V+2z$, is the connection-variable
equation.  The scalar definition $q=K_{\rm phys}-2T$ recovers the trace
equation including the disclosed Appendix-C scalar correction.  Thus these
are the full linearized free Z4c equations in any subsequently chosen
regular physical coordinate gauge.  At this background the gauge changes
leave $H,M,T,z$ invariant, because their background values vanish.

\subsection{Why the reduced solution satisfies the full sourced equations}

Here is a linear proof that avoids assuming a traceful-matter extension of a
published trace-free CEFE constraint theorem.  The initial phase is required
to be a compatible sourced CEFE lift of the physical three-metric and
extrinsic curvature, with the regular rescaled Weyl data.  The harmonic and
scalar-curvature gauge data are chosen consistently.  On any compact subset
of $\Omega>0$, solve the ordinary linear sourced physical Einstein equations
in Lorenz gauge.  Since the physical background is flat, these are
\[
 \widetilde\Box\widehat h_{ab}=-2\Omega^2S_{ab}.
\]
The Lorenz constraint satisfies a homogeneous wave equation by physical
source conservation.  Its initial value is fixed to zero by the four gauge
choices; its initial normal derivative is zero by the sourced Hamiltonian
and momentum equations.  Thus this solution is a genuine physical Ricci
solution on the interior domain of dependence.

Locally put that solution into the prescribed conformal scalar-curvature
and generalized harmonic coordinate gauges.  Infinitesimal conformal changes
obey $\delta_\vartheta R=-6\Box\vartheta-2R_*\vartheta$; infinitesimal
coordinate gauge changes obey the usual vector wave operator
$\Box\xi^a+R_*{}^a{}_b\xi^b$, with a regular first-derivative coupling to
$\vartheta$.  These are ordinary coupled linear waves on each such interior
domain.  Their Cauchy values and first derivatives can be chosen to match
the reduced initial metric and conformal-factor data; the unused normal
gauge derivatives enforce the prescribed gauge sources.  The resulting
genuine rescaled Weyl field satisfies the same regular Bianchi equation with
the same $J$ and the same initial $d$.  Bianchi uniqueness on the interior
domain of dependence therefore identifies it with the independently evolved
field $d$.  The genuine fields satisfy the exact traceful equations above and therefore
the reduced equations with the same initial data.  Uniqueness for that
regular reduced system makes them coincide with its solution.

This interior argument covers every point with $\Omega>0$ at finite time.
The backward background domain of dependence of such a point remains
strictly inside the ball: the glancing radial speed is $O(\Omega^2)$ and
cannot reach $\Omega=0$ in finite backward time.  Thus no unspecified exterior
data enter the comparison.  All sourced CEFE zero quantities vanish in the
interior.  Their regular Sobolev traces then vanish at scri as well.
This proves equivalence for the compatible data class without invoking
a nonexistent arbitrary-trace source theorem.

\subsection{Initial compatibility and what is, and is not, gained}

The mixed initial phase must include a regular curvature lift.  A simple
sufficient way of obtaining $L(0)\in H^4$ is to start with the conformal
three-metric in $H^6$, its conformal extrinsic curvature in $H^5$, and
compatible $d(0)\in H^4$.  Alternatively one can include the compatible
$L(0)\in H^4$ lift explicitly in the phase.  It is incorrect to infer this
order from a bare three-metric in $H^5$, since the Gauss equation has two
spatial derivatives.  The trace and normal derivatives of $L$ are then
fixed by the prescribed scalar curvature, its curl/divergence equations,
and the initial sourced curvature identities.

There is no hidden division in constructing the initial $s$.  If
$K_{\rm phys}=\Omega\bar K+3w$, the physical Hamiltonian identity gives
\[
 s=\frac{\bar\Delta\Omega+\bar K w}{3}
 +\frac{\Omega}{12}
 (R[\bar\gamma]+\bar K^2-\bar K_{ij}\bar K^{ij})
 -\frac{H}{12\Omega}.
\]
The last term is $-\Omega^3h_4/12$ in the present data class.  Initial
$L_{ij},L_{0i}$ can likewise be obtained from the unphysical Gauss--Codazzi
curvature after subtracting $\Omega d$, and $L_{00}$ from its scalar trace;
these operations use derivatives but no inverse conformal factor.

The conclusion of this section is a bounded particular \emph{linear}
physical metric reconstruction in the mixed phase above, including its
constraint-generated Ricci source.  A choice of regular Weyl data is genuine
physical information.  The section does not replace the remaining forced
elliptic-shift/lapse mapping or its initial-data bijection.  The primary
wave-CEFE references
\url{https://arxiv.org/abs/1902.01816} and
\url{https://arxiv.org/abs/2303.07385} explicitly treat trace-free matter;
the traceful equations and the interior-equivalence argument here supply
the additional step needed for this particular prescribed linear source.
